% preamble.tex

% --- Document Structure and Layout ---

\usepackage{pdflscape}        % For landscape pages

% --- Configure fancy headers and footers ---
% \usepackage{fancyhdr}
% \pagestyle{fancy}
% \fancyhf{} % Clear header and footer
% \fancyhead[LE,RO]{\thepage} % Page numbers: Left Even, Right Odd
% \fancyhead[RE]{\nouppercase{\leftmark}} % Chapter title on Even pages
% \fancyhead[LO]{\nouppercase{\rightmark}} % Section title on Odd pages

% --- Paragraph settings ---
% \setlength{\parindent}{1.5em}
% \setlength{\parskip}{0.5em}

% --- Fonts and Encoding ---

\usepackage{fontspec}         % Allows font specification
\usepackage{amsmath}          % For math symbols
\usepackage{luatexja}         % For CJK scripts in LuaLaTeX, load after fontspec

% Main Font
\setmainfont[
  UprightFont = *-Regular,
  ItalicFont = *-Italic,
  ItalicFeatures = { SmallCapsFont = *-Italic },
  SlantedFont = *-Regular,
  SlantedFeatures= { FakeSlant=0.2 },
  BoldFont = *-Bold,
  BoldFeatures = { SmallCapsFont = *-Bold },
  BoldItalicFont = *-BoldItalic,
  BoldItalicFeatures = { SmallCapsFont = *-BoldItalic },
  BoldSlantedFont= *-Bold,
  BoldSlantedFeatures= { FakeSlant=0.2, SmallCapsFont = *-Bold },
  SmallCapsFont = *-Regular,
  SmallCapsFeatures={ RawFeature=+smcp },
  Ligatures=TeX,
  Numbers={OldStyle, Proportional}
]{StixTwoText}

% Math Font
\setmathfont{StixTwoMath.otf}

% Monospace Font
\setmonofont[
  Scale=0.89
]{FiraCode Nerd Font}

% --- Packages for Tables ---

\usepackage{longtable}        % For long tables
\usepackage{array}            % For table column width specification
\usepackage{booktabs}         % For table rules
\usepackage{ragged2e}         % For text alignment (used with \newcolumntype)

% --- Other Packages ---

\usepackage[most]{tcolorbox}  % For colored boxes, load after xcolor to avoid option clash
\usepackage[x11names]{xcolor} % Required for specifying custom colors, load before tcolorbox
\usepackage[version=4]{mhchem}% Formula subscripts using \ce{}

% --- Color Definitions ---

\definecolor{headingblue}{RGB}{23,48,191}
\definecolor{boxtitle}{HTML}{F0F4F8}
\definecolor{boxbody}{HTML}{FBFDFF}
\definecolor{mainboxframe}{HTML}{F0F4F8}
\definecolor{subboxframe}{HTML}{F0F4F8}

% --- Boxes --- (WORKS!)

% Define the floating environment
\newfloat{scholarlybox}{tbp}{box}[chapter]
\floatname{scholarlybox}{Box}

% Configure counter
\makeatletter
\@addtoreset{scholarlybox}{chapter}
\renewcommand{\thescholarlybox}{\thechapter.\arabic{scholarlybox}}
\makeatother

% Box construction command
\newcommand{\boxcontent}[2]{%
    \begin{scholarlybox}
        \refstepcounter{scholarlybox}
        \setlength{\fboxrule}{0.5pt}%
        \setlength{\fboxsep}{0pt}%
        \fbox{%
            \setlength{\fboxsep}{10pt}%
            \colorbox{gray!5}{%
                \begin{minipage}{\dimexpr\linewidth-2\fboxsep-2\fboxrule\relax}
                    {\textsc{Box~\thescholarlybox\quad#1}\par}
                    \vspace{6pt}
                    {\RaggedRight#2}
                \end{minipage}%
            }%
        }%
        \label{box:#1}%
    \end{scholarlybox}
}

% --- sansblock Environment ---

\usepackage{sourcesanspro}    % Load Source Sans Pro
\setsansfont{Source Sans Pro} % Set it as the sans-serif font
\newenvironment{sansblock}[1]
    {\small\sffamily\raggedright{\scshape #1}\ } % Ensure small caps for the title
  {} % End environment: no special commands needed

% --- Custom Column Type (using ragged2e) ---

\newcolumntype{R}[1]{>{\RaggedRight}p{#1}}

% ---  Margin notes ---

\usepackage{marginnote}
\renewcommand*{\marginfont}{\footnotesize\itshape} % Style for margin notes

% --- Epigraph ---

\usepackage{epigraph}
\setlength\epigraphwidth{.9\textwidth}
\newenvironment{quotepara}
  {\setlength{\parindent}{1em}\setlength{\parskip}{0.6ex}\itshape\raggedright}
  {}
\renewcommand{\textflush}{quotepara}

% --- Small Caps ---

\newcommand{\flatcaps}[1]{\textsc{\MakeLowercase{#1}}}

% --- Lists ---

% General settings for all lists
\usepackage{enumitem}
\setlist{nosep} % Removes vertical spacing between items and lists

% % Specific settings for level-1 lists
\setlist[1]{
  labelindent=0pt,   % Set label indent to 0
  leftmargin=*,      % Use automatic calculation for left margin based on label width
  itemindent=0pt,    % No extra indentation for item text
  labelsep=0.5em,   % Space between bullet and text
  align=left,       % Ensure label is aligned to the left
  % label=\textbullet % Explicitly set the label to a bullet (optional, but good for clarity)
}

% --- Custom Chapter/Section Styles ---

\usepackage[compact]{titlesec}  % Allows creating custom chapter styles
\titleformat{\chapter}[display]
  {\fontsize{60}{62}\bfseries}
  {\thechapter}
  {0pt}
  {\huge\noindent}
\titlespacing*{\chapter}{0pt}{0pt}{40pt}

\titleformat{\section}
  {\normalsize\normalfont}
  {\thesection}
  {0.6em}
  {\flatcaps}
\titlespacing*{\section}{0pt}{1\baselineskip}{1\baselineskip}

\titleformat{\subsection}[block]
  {\normalsize\normalfont} % defines the font size and style for the entire subsection heading, including both the number and the title
  {\thesubsection} % defines the format of the subsection number
  {1em} % the horizontal space between the subsection number and the title
  {\itshape} % defines the format of the subsection title itself
\titlespacing*{\subsection}{0pt}{1\baselineskip}{1\baselineskip}

\titleformat{\subsubsection}[runin]
  {\normalsize\normalfont} % defines the font size and style for the entire subsection heading, including both the number and the title
  {\thesubsubsection} % defines the format of the subsection number
  {1em} % the horizontal space between the subsection number and the title
  {\itshape}[.] % defines the format of the subsection title itself
\titlespacing*{\subsubsection}{0pt}{1\baselineskip}{1\baselineskip}

\titleformat{\paragraph}[runin]
  {\flatcaps}
  {\theparagraph}
  {0pt}
  {}

% --- Footnotes ---

\usepackage[norule,ragged,hang]{footmisc}  % Load footmisc with ragged option
\renewcommand{\footnotelayout}{\RaggedRight\footnotesize} % Typeset footnotes in \RaggedRight
\setlength{\footnotemargin}{1.5em}   % Adjust space between number and text
\makeatletter
\renewcommand{\@makefntext}[1]{%
    \parindent 1em%                    Set parindent for footnote text
    \noindent
    \hb@xt@ 1.8em{%                   Set hanging indent for footnote text
        \hss\@thefnmark.%
    }
    \RaggedRight #1%                  Typeset footnote text ragged right
}
\makeatother

% --- Captions ---

\usepackage{caption}
\captionsetup{
  font={small},
  labelfont={bf},
  textfont={},
  width=0.9\textwidth,
  justification=justified,
  labelformat=default,
  labelsep=period,
  format=plain
}
\renewcommand{\captionlabelfont}{\bfseries\scshape}

% --- Hyperlinks ---

\usepackage{hyperref}           % Load after most other packages, but before cleveref
\hypersetup{
  colorlinks = true,
  linkcolor = {headingblue},
  citecolor = {headingblue},
  urlcolor = {headingblue}
}

% --- Miscellaneous ---

\usepackage{iftex}             % Detects the engine used
\usepackage{microtype}         % Improves typography
\AddToHook{env/Highlighting/begin}{\small} % Set the code chunk font size globally
\usepackage{etoolbox}
%% Use lining fonts
\newcommand\lining{\addfontfeatures{Numbers={Monospaced, Lining}}}
\AtBeginEnvironment{tabular}{\lining} % In tables
\renewcommand{\theequation}{ {\lining\arabic{equation}}} % For equation numbers

% --- Index (if needed) ---

\usepackage{makeidx}
\makeindex

% --- Other Typography Settings ---

\frenchspacing                % Single space after periods
\tolerance=400                % Default is 200; higher values allow more relaxed line-breaking.
\emergencystretch=3em         % Adds additional space to help line-breaking.
\hyphenpenalty=20             % Default is 50; lower values encourage hyphenation.

% --- Microtype Settings (adjust only if needed) ---

% \microtypesetup{
%    tracking = true,
%    factor = 1100,
%    stretch = 15,
%    shrink = 15
% }
